%
% Übersicht über den Algorithmus
%
\subsection{Übersicht über den Algorithmus
\label{buch:intalg:risch:uebersicht}}
Bevor wir die notwendigerweise unterschiedlichen Vorgehensweisen 
für logarithmische, exponentielle und algebraische Erweiterungen
im Detail betrachten, wollen wir in diesem Abschnitt die für alle
Fälle gemeinsame prinzipielle Vorgehensweise
des Algorithmus skizzieren.

%
% Die Aufgabenstellung
%
\subsubsection{Die Aufgabenstellung}
Es ist ein Algorithmus zu entwickeln, der zu einer elementaren Funktion
eine elementare Stammfunktion findet, sofern eine solche existiert.
Falls keine elementare Stammfunktion existiert, soll der Algorithmus
einen Beweis für diese Tatsache konstruieren.
Um dieses Ziel zu erreichen, muss zunächst klargestellt werden, welche
Funktionen als Integranden in Frage kommen sollen.

Zu den Funktionen, die als Stammfunktionen zugelassen werden sollen,
sollen mindestens alle rationalen Funktionen mit rationalen Koeffizienten
gehören.
Wir gehen also von einem Funktionenkörper aus, der eine Erweiterung
von $\mathbb{Q}$ um ein transzendentes Element $x$ ist, welches die
Eigenschaft $Dx=1$ hat.
Dieser Körper wird auch mit $\mathbb{Q}(x)$ bezeichnet.

Ein typisches Computeralgebrasystem wird noch weitere Konstanten
als Koeffizienten der rationalen Funktionen zulassen, zum Beispiel
Naturkonstanten wie $g$ oder $c$, denen man zwar auch einen numerischen
Wert zuordnen kann, aber man interessiert sich durchaus auch dafür, 
wie solche Konstanten in die Stammfunktion eingehen können.
Solche Konstanten können als transzendente Erweiterungen des Körpers
$\mathbb{Q}$ angesehen werden, die zusätzlich die Eigenschaft haben,
dass ihre Ableitung verschwindet.
Weiter ist es möglich, dass zusätzliche mathematische Konstanten
benötigt werden.
Diese können sowohl transzendente Erweiterungen sein wie $\pi$, $e$,
$\zeta(2)=\pi^2/6$ oder $\zeta(3)$, oder sie können algebraisch sein wie
$\!\sqrt{2}$ oder $\sqrt[3]{7}$.
Die algebraische Abhängigkeit dieser Konstanten untereinander wird
nicht weiter untersucht und hat auch keinen Einfluss auf das
Integrationsproblem.

Innerhalb des Konstantenkörpers kann es durchaus passieren, dass eine
transzendente Konstante in einer algebraischen Erweiterung von
$\mathbb{Q}$ liegt.
Dies passiert zum Beispiel für die transzendente Konstate $\pi$,
$\mathbb{Q}\subset \mathbb{Q}(\pi)$ ist eine transzendente Erweiterung.
Sie enthält auch $\zeta(2) = \pi^2/6$.
Andererseits ist $\mathbb{Q}(\zeta(2)) \subset \mathbb{Q}(\zeta(2),\pi)$
eine algebraische Erweiterung, da $\pi$ eine Nullstelle des Polynoms
\[
p(X)
=
X^2 - 6\zeta(2)
=
X^2 - \pi^2
\qquad
\Rightarrow
\qquad
p(\pi)=0
\]
ist.
Während $\mathbb{Q}(\pi)$ auch $\zeta(2)$ bereits enthält, braucht
es eine algebraische Erweiterung von $\mathbb{Q}(\zeta(2))$, um $\pi$
zu finden.
Im Folgenden nehmen wir an, dass alle nötigen Konstanten von Anfang
an dem Körper $\mathbb{Q}$ hinzugefügt worden sind und den Konstantenkörper
$K$ bilden.

Es ist auch denkbar, dass während der Bestimmung der Stammfunktion weitere
Konstanten benötig werden. 
Zum Beispiel kann es nötig sein, ein Polynom zu faktorisieren.
Die dabei gefundenen Nullstellen sind algebraische Konstanten über
dem Körper $K$.
Dazu gehört insbesondere auch das Element $i$ mit $i^2=-1$, welches
zum Beispiel für die Konstruktion der trigonometrischen Funktionen
benötigt wird.
Auch diese Konstanten sollen von Anfang an in $K$ enthalten sein, so
dass es später nicht mehr nötig ist, den Konstantenkörper zu erweitern.

Die Konstruktion des differentiellen Funktionenkörpers beginnt daher
mit dem trans\-zen\-den\-ten Element $x$ und dem Körper $F=K(x)$ der rationalen
Funktionen in $x$ mit Koeffizienten in $K$.
Diesem Körper werden jetzt weiter drei Arten von Funktionen hinzugefügt:
\begin{itemize}
\item
Logarithmische transzendente Erweiterungen: Dies sind Funktionen $\vartheta=\log u$,
die durch die Eigenschaft $D\vartheta= Du/u$ bis auf eine Konstante
definiert sind.
\item
Exponentielle transzendente Erweiterungen: Dies sind Funktionen $\vartheta=e^u$,
die durch die Eigenschaft $D\vartheta = \vartheta\cdot Du$ bis auf einen
skalaren Faktor definiert sind.
\item
Algebraische Elemente: Dies sind Funktionen $\vartheta$, die eine algebraische
Gleichung
\[
a_n\vartheta^n + a_{n-1}\vartheta^{n-1} + \dots + a_1\vartheta + a_0 = 0
\]
erfüllen.
\end{itemize}
Um einen zulässigen Integranden oder eine Stammfunktion zu konstruieren,
kann es notwendig sein, mehrere solche Erweiterungen hinzufügen, die auch
von verschiedener Art sein können.
Wir haben also mit einer Folge von Erweiterungen
\[
F
\subset
\underbrace{
F(\vartheta_1)
}_{\displaystyle = F_1}
\subset
\underbrace{
F(\vartheta_1,\vartheta_2)
}_{\displaystyle = F_2}
\subset
\dots
\subset
\underbrace{
F(\vartheta_1,\vartheta_2,\dots,\vartheta_n)
}_{\displaystyle=F_n}
\]
zu tun.
Insbesondere ist auch
\[
F_n = F_{n-1}(\vartheta_n)=F(\vartheta_1,\dots,\vartheta_{n-1})(\vartheta_n).
\]
Dieser Aufbau des Funktionenkörpers durch wiederholte Adjunktion von 
logarithmischen, exponentiellen oder algebraischen Elementen war auch
bereits der Rahmen, in dem das Liouville-Prinzip formuliert und im
Abschnitt~\ref{buch:intalg:liouville:subsection:allgemein}
bewiesen wurde.

Mit den Differentialkörpern $F_k$ kann jetzt die Aufgabenstellung 
genauer formuliert werden.
Zu einer Funktion $f\in F_n$ soll der Algorithmus eine Stammfunktion
$g\in F_n$ finden, also eine Funktion $Dg=f$.
Falls keine solche Funktion existiert, soll er einen Beweis dafür
konstruieren.

%
% Der polynomielle und der rationale Teil
%
\subsubsection{Der polynomielle und der rationale Teil}
Wir betrachten eine Funktion $f\in F_n = F_{n-1}(\vartheta_n)$ und
versuchen die Frage zu beantworten, ob $f$ eine Stammfunktion in
$F_n$ hat.
Dies bedeutet, dass $f$ als eine rationale Funktion der Form
\begin{equation}
f
= 
\frac{p(\vartheta_n)}{q(\vartheta_n)},
\qquad
p(\vartheta_n), q(\vartheta_n)\in F_{n-1}[\vartheta_n]
\label{buch:intalg:risch:uebersicht:eqn:integrand}
\end{equation}
geschrieben werden kann.
Im Gegensatz zur Situation in Kapitel~\ref{chapter:ratint}, wo rationale
Funktionen mit konstanten Koeffizienten untersucht wurden, haben
die Polynome $p(\vartheta_n)$ und $q(\vartheta_n)$ Koeffizienten
in $F_{n-1}$, sind also selbst Funktionen.

Auf die Polynome im Integranden
\eqref{buch:intalg:risch:uebersicht:eqn:integrand}
kann der Polynomdivisionsalgorithmus angewendet werden, was auf die
Zerlegung
\[
\frac{p(\vartheta_n)}{q(\vartheta_n)}
=
s(\vartheta_n)
+
\frac{r(\vartheta_n)}{q(\vartheta_n)},
\qquad
s(\vartheta_n), q(\vartheta_n) \in F_{n-1}[\vartheta_n],
\]
führt.
Das Polynom $r(\vartheta_n)$ hat Grad
$\deg r(\vartheta_n) < \deg q(\vartheta_n)$.
Wir nennen $s(\vartheta_n)$ den polynomiellen Teil und den Bruch
den rationalen Teil.

Das gesuchte Integral kann jetzt in die zwei Integrale
\[
\int f
=
\int
\frac{p(\vartheta_n)}{q(\vartheta_n)}
=
\underbrace{
\int s(\vartheta_n)
}_{\displaystyle\text{polynomieller Teil\strut}}
+
\underbrace{
\int
\smash{\frac{r(\vartheta_n)}{q(\vartheta_n)}}.
}_{\displaystyle\text{rationaler Teil\strut}}
\]
zerlegt werden.
Im Falle rationaler Funktionen war der polynomielle Teil sehr einfach
auszuführen, da $\int z^n=z^{n+1}/(n+1)$.
Im vorliegenden, allgemeineren Fall wird eine neue Methode nötig.
Schon das einfache Polynom $s(\vartheta_n) = \vartheta_n$ führt auf
das nichttriviale Integral $\int s(\vartheta_n) = \int\vartheta_n$.
Wegen der sehr unterschiedlichen Eigenschaften der verschiedenen Arten
von Erweiterungen wird dies jeweils unterschiedliche Vorgehensweisen
erfordern, die in den nachfolgenen Abschnitten für jede Erweiterung
separat entwickelt werden müssen.
Die in Kapitel~\ref{chapter:ratint} mit einer Randbemerkung erledigte
Behandlung des polynomiellen Teils muss daher nochmals neu aufgerollt
werden und wird zu zusätzlichen Einschränkungen für die Existenz
einer Stammfunktion führen.

Für den rationalen Teil können wir ebenfalls versuchen, die
bereits bekannten Techniken anzuwenden.
Wir können eine quadratfreie Faktorisierung des Nenners $q(\vartheta_n)$
in Faktoren $b_k(\vartheta_n)\in F_{n-1}[\vartheta_n]$ vornehmen und
damit den Bruch in Partialbrüche aufteilen.
Weiter können wir die Methode von Hermite auf die Partialbrüche 
anwenden und damit den Nennerexponenten auf $1$ reduzieren.
Schliesslich können wir die Methode von Rothstein-Trager anwenden,
um die Koeffizienten $c_i$ und die Funktion $v_i\in F_{n-1}$ für
die Liouville-Darstellung der Stammfunktion zu finden.
Dabei wird sich allerdings eine neue Einschränkung ergeben, die darauf
beruht, dass die $c_i$, die die Methode von Rothstein-Trager findet,
Konstanten sein müssen.
Da die Koeffizienten der Zähler- und Nennerpolynome nicht konstant sein
müssen, könnten sich auch nicht konstante $c_i$ ergeben, was dann
beweist, dass die Stammfunktion nicht existieren kann.

Der Integrationsalgorithmus in $F_n$ folgt also ziemlich genau der
Vorgehensweise des Integrationsalgorithmus in $K(z)$.
Um die Notation zu vereinfachen, ist es daher angebracht, $F=F_{n-1}$
und $\vartheta=\vartheta_{n}$ zu setzen.
Zu $f\in F(\vartheta)=F_{n-1}(\vartheta_n)=F_n$ ist dann eine Stammfunktion
in $F(\vartheta)$ gesucht.
Im Gegensatz zur Situation in Kapitel~\ref{chapter:ratint} ist
aber nicht garantiert, dass die Stammfunktion tatsächlich existiert.

%
% Die Methode von Hermite
%
\subsubsection{Die Methode von Hermite}
Die Methode von Hermite verwendet eine quadratfreie Faktorisierung 
\[
q(\vartheta)
=
b_1(\vartheta) b_2(\vartheta)^2 \cdots b_m(\vartheta)^m,
\]
in der jeder Faktor $b_k(\vartheta)$ quadratfrei ist.
Zu dieser Faktorisierung gibt es immer eine Partialbruchzerlegung
\begin{equation}
\frac{r(\vartheta)}{q(\vartheta)}
=
\sum_{i=1}^m \sum_{k=1}^i \frac{a_{ik}(\vartheta)}{b_i(\vartheta)^k}
\end{equation}
mit $a_{ik}(\vartheta)\in F[\vartheta]$ und
$\deg a_{ik}(\vartheta)<\deg b_i(\vartheta)$.
Das Integral des rationalen Teils ist damit
\begin{equation}
\int
\frac{r(\vartheta)}{q(\vartheta)}
=
\sum_{i=1}^m \sum_{k=1}^i
\int \frac{a_{ik}(\vartheta)}{b_i(\vartheta)^k}.
\end{equation}
Es ist also nötig, Integrale der Form
\begin{equation}
\int \frac{a(\vartheta)}{b(\vartheta)^k}
\qquad
\text{mit $
a(\vartheta), b(\vartheta)\in F[\vartheta]
$}
\label{buch:intalg:risch:eqn:ratpart}
\end{equation}
zu bestimmen, mit dem quadratfreien Nennerpolynom
$b(\vartheta)$ mit $\deg a(\vartheta) < \deg b(\vartheta)$.
Die Methode von Hermite versucht, die
Integrale~\eqref{buch:intalg:risch:eqn:ratpart}
zu vereinfachen, indem der Nennerexponent $k$ reduziert wird.

Da $b(\vartheta)$ quadratfrei ist, sind $b(\vartheta)$ und $Db(\vartheta)$
in den meisten Fällen teilerfremd.
Die Fälle von Faktoren $\vartheta$ des Nennerpolynoms, in denen dies nicht
zutrifft, hängen von der Art der Erweiterung $\vartheta$ ab und werden
später separat diskutiert.
Teilerfremdheit hat zur Folge, dass es Polynome $s(\vartheta)$ und
$t(\vartheta)$ gibt mit
\[
a(\vartheta)
=
s(\vartheta) b(\vartheta)
+
t(\vartheta) Db(\vartheta).
\]
Damit wird das Integral zu
\begin{align*}
\int
\frac{a(\vartheta)}{b(\vartheta)^k}
&=
\int
\frac{s(\vartheta)}{b(\vartheta)^{k-1}}
+
\int t(\vartheta)\frac{Db(\vartheta)}{b(\vartheta)^k}
\intertext{und mit partieller Integration wird daraus}
&=
\int
\frac{s(\vartheta)-Dt(\vartheta)/(k-1)}{b(\vartheta)^{k-1}}
+
\frac{t(\vartheta)/(k-1)}{b(\vartheta)^{k-1}}.
\end{align*}
Der Zähler des ersten Integrals hat möglicherweise grösseren Grad
als $b(\vartheta)$.
Die Polynomdivision macht daraus einen zusätzlichen Beitrag zum polynomiellen
Teil.
Es bleibt das Problem, Integrale der Form~\eqref{buch:intalg:risch:eqn:ratpart}
mit $k-1$ anstelle von $k$ zu bestimmen.

Durch wiederholte Anwendung dieser Methode ist das Problem damit reduziert
auf Integrale der Form
\[
\int\frac{a(\vartheta)}{b(\vartheta)},
\qquad
a(\vartheta), b(\vartheta) \in F[\vartheta],
\]
wobei $b(\vartheta)$ quadratfrei ist.

%
% b(\vartheta) und Db(\vartheta)
%
\subsubsection{$b(\vartheta)$ und $Db(\vartheta)$}
Der entscheidende Schritt für die Durchführbarkeit der Methode von Hermite 
war die Garantie, dass
\[
a(\vartheta)
=
s(\vartheta) b(\vartheta)
+
t(\vartheta) Db(\vartheta)
\]
geschrieben werden kann.
Falls $b(\vartheta)$ und $Db(\vartheta)$ teilerfremd sind, dann gilt
dies nach dem erweiterten euklidischen Algorithmus tatsächlich.
Für ein quadratfreies Polynom in $u(X)\in K[X]$ wissen wir auch bereits,
dass $u(X)$ und das abgeleitete Polynom $u'(X)$
(siehe Definition~\ref{buch:intalg:elementar:def:abgeleitetespolynom})
teilerfremd sind.
Bei der Ableitung $Db(\vartheta)$ werden aber nicht nur die Potenzen
$\vartheta^k$ abgeleitet, sondern auch die Koeffizienten des Polynoms.

Man beachte, dass $Db(\vartheta)$ und $b'(\vartheta)$ verschieden sind.
Nach Definition~\ref{buch:intalg:elementar:def:abgeleitetespolynom} gilt
\begin{align*}
Db(\vartheta)
&=
\sum_{k=0}^m
D(b_k\vartheta^k)
\\
&=
\sum_{k=0}^m
\Bigl(
Db_k\vartheta^k
+
b_kk D\vartheta\cdot \vartheta^{k-1}
\Bigr)
\\
&=
\sum_{k=1}^m (Db_k)\vartheta^k
+
b'(\vartheta) D\vartheta.
\end{align*}
Da $D\vartheta$ auf der rechten Seite vorkommt, ist zunächst nicht einmal
klar, dass die rechte Seite auch wieder ein Polynom in $\vartheta$ mit
Koeffizienten in $F$ ist.
Dies muss später für jede Art von Erweiterung einzeln nachgeprüft werden.

Aus der Quadratfreiheit von $b(\vartheta)$ als Polynom in $F[\vartheta]$
folgt nur, dass $b(\vartheta)$ und $b'(\vartheta)$ teilerfremd sind.
Damit die Methode von Hermite durchführbar wird, muss also untersucht
werden, ob aus der Teilerfreiheit von $b(\vartheta)$ und $b'(\vartheta)$
auch folgt, dass $b(\vartheta)$ und $Db(\vartheta)$ teilerfremd sind.
Dies wird für logarithmische Erweiterungen in
Lemma~\ref{buch:intalg:risch:log:lemma:bDb} und für exponentielle
Erweiterungen in
Lemma~\ref{buch:intalg:risch:exp:lemma:bDb} durchgeführt.

%
% Standardform der Stammfunktion
%
\subsubsection{Standardform der Stammfunktion}
Die Methode von Rothstein-Trager von
Satz~\ref{buch:ratint:rothsteintrager:satz:rothsteintrager}
bestimmt die Stammfunktion einer rationalen Funktion mit quadratfreiem
Nenner und einem Zähler kleineren Grades und integriert damit den
rationalen Teil eines rationalen Integranden.
Die Methode soll auf Funktionen 
\begin{equation}
\frac{a(\vartheta)}{b(\vartheta)}
\qquad \text{mit}\quad \left\{
\begin{array}{l}
a(\vartheta),b(\vartheta)\in F[\vartheta],\\
\text{$a(\vartheta)$ und $b(\vartheta)$ teilerfremd,}\\
\text{$b(\vartheta)$ quadratfrei und normiert,}\\
\deg a(\vartheta) < \deg b(\vartheta)\\
\end{array}
\right.
\label{buch:intalg:risch:eqn:a(theta)/b(theta)}
\end{equation}
ausgedehnt werden.
Möglicherweise müssen die zulässigen Nenner $b(\vartheta)$ weiter
eingeschränkt werden.
Zum Beispiel wird sich herausstellen, dass bei einer exponentiellen
Erweiterung $\vartheta$ für den rationalen Teil eines
Integrals nur Nenner $b(\vartheta)$ zulässig sind, die nicht durch
$\vartheta$ teilbar sind.
Wir sagen, dass \eqref{buch:intalg:risch:eqn:a(theta)/b(theta)}
ein \emph{rationaler Teil} für die Körpererweiterung $F\subset F(\vartheta)$
ist, wenn $b(\vartheta)$ ein zulässiger Nenner ist.

Nach dem Liouville-Prinzip muss die Stammfunktion, falls sie elementar ist,
die Form
\begin{equation}
\int
\frac{a(\vartheta)}{b(\vartheta)}
=
v_0(\vartheta)
+
\sum_k c_k\log v_k(\vartheta)
\label{buch:intalg:risch:rt:eqn:liouvilleint}
\end{equation}
haben.
Für den Integranden muss entsprechend gelten, dass
\begin{equation}
\frac{a(\vartheta)}{b(\vartheta)}
=
Dv_0(\vartheta)
+
\sum_k c_k\frac{D v_k(\vartheta)}{v_k(\vartheta)}
\label{buch:intalg:risch:rt:eqn:liouvillediff}
\end{equation}
ist.
Darin sind die $c_k$ Konstanten und $v_k(\vartheta)\in F(\vartheta)$.
Die Methode von Rothstein-Trager findet die $c_k$ und die $v_k(\vartheta)$.

Die Darstellung
\eqref{buch:intalg:risch:rt:eqn:liouvilleint}
der Stammfunktion ist nicht die einzig mögliche.
Auch bei der Integration rationaler Funktionen waren Umformungen
verwendet worden, um der Stammfunktion genau die Form zu geben,
die die Methode von Rothstein-Trager später fand.
Im Folgenden sollen ähnliche Vereinheitlichungen für die Darstellung
\eqref{buch:intalg:risch:rt:eqn:liouvilleint}
vorgenommen werden.

In der Darstellung \eqref{buch:intalg:risch:rt:eqn:liouvilleint}
der Stammfunktion sind die $v_k(\vartheta) \in F(\vartheta)$ rationale
Funktionen von $\vartheta$.
Im Beweis des
Satzes~\ref{buch:ratint:rothsteintrager:satz:rothsteintrager}
über die Stammfunktion rationaler Funktionen
wurde das Wissen aus der Methode von Bernoulli verwendet, dass die
\index{Bernoulli, Methode von}%
\index{Methode von Bernoulli}%
\index{Bernoulli, Johann}%
$v_k$ auch als normierte Polynome gewählt werden können.
Dies trifft aber auch im vorliegenden Fall zu.
Sind nämlich die $v_k(\vartheta)$ Brüche
\[
v_k(\vartheta)
=
\frac{z_k(\vartheta)}{w_k(\vartheta)}
\]
von Zähler $z_k(\vartheta)$ und Nenner $w_k(\vartheta)$
mit $z_k(\vartheta), w_k(\vartheta)\in F[\vartheta]$, dann können
wir \eqref{buch:intalg:risch:rt:eqn:liouvilleint} in der Form
\begin{equation*}
\int
\frac{a(\vartheta)}{b(\vartheta)}
=
v_0(\vartheta)
+
\sum_k c_k\log v_k(\vartheta)
=
v_0(\vartheta)
+
\sum_k c_k\log z_k(\vartheta)
-
\sum_k c_k\log w_k(\vartheta)
\end{equation*}
schreiben.
Die Argumente der Logarithmen sind jetzt sogar Polynome $\vartheta$.
Dies zeigt, dass man annehmen darf, dass in
\eqref{buch:intalg:risch:rt:eqn:liouvilleint}
$v_k(\vartheta)\in F[\vartheta]$ gilt.

Wir möchten zeigen, dass man annehmen darf, dass $v_k(\vartheta)$
normiert ist.
Dazu schreiben wir
$v_k(\vartheta) = v_{\deg v_0(\vartheta)} \bar{v}_k(\vartheta)$, worin
$\bar{v}_k(\vartheta)$ ein normiertes Polynom ist.
Dann ist
\[
c_k\log v_k(\vartheta)
=
c_k\log v_{\deg v_k(\vartheta)}
+
c_k\log \bar{v}_k(\vartheta).
\]
Indem wir den Ausdruck $c_k\log v_{\deg v_k(\vartheta)}\in F$ zu
$v_0(\vartheta)$ hinzufügen, können wir annehmen, das $v_k(\vartheta)$
\emph{normiert} ist.

Falls mehrere der Koeffizienten $c_k$ gleich sind, dann können wir 
die zugehörigen Funktionen $v_k$ zusammenfassen.
Zu gegebenem $c$ schreiben wir daher
\[
v(\vartheta) = \prod_{c_k=c} v_k(\vartheta).
\]
Dann ist
\begin{align*}
c \log v(\vartheta)
&=
c
\log\biggl(\prod_{c_k=c} v_k(\vartheta)\biggr)
=
\sum_{c_k=c} c_k \log v_k(\vartheta)
\intertext{mit der Ableitung}
c \frac{Dv(\vartheta)}{v(\vartheta)}
&=
D(c\log v(\vartheta))
=
\sum_{c_k=c} c_k \frac{Dv_k(\vartheta)}{v_k(\vartheta)}.
\end{align*}
Man darf also immer annehmen, dass die $c_k$ in der Liouville-Form der
Stammfunktion alle \emph{verschieden} sind.

Falls zwei Polynome $v_k(\vartheta)$ und $v_l(\vartheta)$ einen nichttrivialen
gemeinsamen Faktor haben, dann können wir mit
$g(\vartheta)=\operatorname{ggT}(v_k(\vartheta),v_l(\vartheta))$
\[
v_k(\vartheta) = w_k(\vartheta) g(\vartheta)
\qquad\text{und}\qquad
v_l(\vartheta) = w_l(\vartheta) g(\vartheta)
\]
schreiben.
Es gilt dann
\[
c_i\log v_i(\vartheta) + c_j\log v_j(\vartheta)
=
c_i\log w_i(\vartheta) + c_j \log w_j(\vartheta) + (c_i+c_j)\log g(\vartheta).
\]
Die Argumente der Logarithmen auf der rechten Seite haben keine 
gemeinsamen Faktoren.
Man darf also annehmen, dass die $v_k(\vartheta)$ \emph{teilerfremd} sind.

Falls $v_k(\vartheta)$ mehrfache Faktoren hat, können diese wie folgt
eliminiert werden.
In der quadratfreien Faktorisierung
\[
v_k(\vartheta) = p_0 p_1(\vartheta) p_2(\vartheta)^2 \dots p_l(\vartheta)^l
\]
von $v_k(\vartheta)$ sind alle Faktoren $p_i(\vartheta)\in F[\vartheta]$
quadratfrei.
Der Logarithmus ist
\[
\log v_k(\vartheta)
=
\log p_0
+
\log p_1(\vartheta)
+
2 \log p_2(\vartheta)
+
\dots
+
l\log p_l(\vartheta)
\]
ein Summe mit lauter quadratfreien Argumenten des Logarithmus.
Man darf also annehmen, dass die $v_k(\vartheta)$ sogar \emph{quadratfrei} sind.

Der Term $v_0(\vartheta)$ kam in der Methode von Rothstein-Trager für
rationale Funktionen nicht vor, da diese nur den logarithmischen Teil
geliefert hat.
Im vorliegenden Fall ist die Sachlage etwas komplizierter.
Sie wird geklärt durch die folgenden zwei Lemmata, sie später
in den Abschnitten~\ref{buch:intalg:risch:logarithmisch}
und \ref{buch:intalg:risch:exponentiell}
für logarithmische bzw.~exponentielle Erweiterungen bewiesen.

\begin{lemma}
\label{buch:intalg:risch:lemma:v0log}
Ist $\vartheta$ eine logarithmische Körpererweiterung von $F$, dann ist
in der Stammfunktion 
\eqref{buch:intalg:risch:rt:eqn:liouvilleint}
eines rationalen Teils immer $v_0(\vartheta)=0$.
\end{lemma}

\begin{lemma}
\label{buch:intalg:risch:lemma:v0exp}
Ist $\vartheta$ eine exponentielle Körpererweiterung von $F$ mit
$D\vartheta = \vartheta Du$, dann ist in der Stammfunktion 
\eqref{buch:intalg:risch:rt:eqn:liouvilleint}
eines rationalen Teils immer $v_0(\vartheta)\in F$.
\end{lemma}

Mit diesen zwei Lemmata ausgerüstet können wir jetzt wieder zeigen,
dass die $v_k(\vartheta)$ so gewählt werden können, dass $b(\vartheta)$
das Produkt der $v_k(\vartheta)$ ist.
Zu diesem Zweck nehmen wir an, dass
\begin{equation}
\frac{a(\vartheta)}{b(\vartheta)}
=
v_0 + \sum_{i=1}^m c_k \frac{Dv_k(\vartheta)}{v_k(\vartheta)},
\label{buch:intalg:risch:uebersicht:eqn:stammfunktion}
\end{equation}
und
verwenden wir die Notation
\[
u_k(\vartheta)
=
\prod_{\scriptstyle i=1\atop\scriptstyle i\ne k} v_i(\vartheta)
\qquad\text{und}\qquad
v(\vartheta)
=
\prod_{i=1} v_i(\vartheta).
\]
Wir wollen zeigen, dass $b(\vartheta)=v(\vartheta)$.

Wir multiplizieren \eqref{buch:intalg:risch:uebersicht:eqn:stammfunktion}
mit $b(\vartheta)$ und $v(\vartheta)$ und erhalten
\[
a(\vartheta) v(\vartheta)
=
v_0 b(\vartheta) \prod_{i=1}^m v_i(\vartheta)
+
\sum_{i=1}^m Dv_i(\vartheta) u_i(\vartheta) b(\vartheta).
\]
$b(\vartheta)$ teilt jeden Term auf der rechten Seite, ist aber teilerfremd
zu $a(\vartheta)$, also muss $b(\vartheta)\mid a(\vartheta)$ teilen.
Andererseits teilt $v_k(\vartheta)$ die linke Seite, den ersten Term
auf der rechten Seite und alle Terme der Summe ausser dem Term mit $i=k$.
Daher muss $v_k(\vartheta)$ auch letzteren Term
\[
c_k Dv_k(\vartheta) u_k(\vartheta) b(\vartheta)
\]
teilen.

Für logarithmische Erweiterungen folgt jetzt bereits aus der Quadratfreiheit
von $v_k(\vartheta)$, dass $v_k(\vartheta) \mid b(\vartheta)$.
Für exponentielle Erweiterungen könnte $v_k(\vartheta)$ auch ein
Teiler von $Dv_k(\vartheta)$ sein.
Dann müsste aber $v_k(\vartheta)$ ein Monom sein, was nicht zulässig
ist, weil der Faktor $\vartheta$ von $b(\vartheta)$ für exponentielle
Erweiterungen ausgeschlossen wurde.
Da jeder Faktor $v_k(\vartheta)$ ein Teiler ist, gilt auch
$v(\vartheta)\mid b(\vartheta)$.
Da sowohl $b(\vartheta)$ wie auch $v(\vartheta)$ normiert sind, folgt
$b(\vartheta) = v(\vartheta)$.

Wir fassen die gefundenen Standardisierungen der Stammfunktion im folgenden
Satz zusammen.

\begin{satz}[Standardform]
\label{buch:intalg:risch:satz:standardform}
\index{Satz!über die Standardform}%
Sei $\vartheta$ eine logarithmische oder exponentielle Erweiterung.
Falls der rationale Teil $a(\vartheta)/b(\vartheta)$ eine elementare
Stammfunktion hat, dann gibt es $v_0\in F$ und normierte, quadratfreie,
paarweise teilerfremde Polynome $v_k(\vartheta)\in F[\vartheta]$ mit
\begin{equation}
\frac{a(\vartheta)}{b(\vartheta)}
=
v_0 + \sum_{k=1}^m \frac{Dv_k(\vartheta)}{v_k(\vartheta)}
\qquad\text{und}\qquad
b(\vartheta)
=
\prod_{k=1}^m v_k(\vartheta).
\label{buch:intalg:risch:satz:eqn:standardform}
\end{equation}
Für eine logarithmische Erweiterung ist $v_0=0$.
\end{satz}

%
% Die Methode von Rothstein-Trager
%
\subsubsection{Die Methode von Rothstein-Trager}
Um die Methode von Rothstein-Trager auf die vorliegende Situation
zu übertragen, übersetzen wir das Resultat des
Satzes~\ref{buch:ratint:rothsteintrager:satz:rothsteintrager}.
Gesucht sind die verschiedenen Elemente $c_k$, für die 
\[
a(\vartheta) - c_k \cdot Db(\vartheta)
\qquad\text{und}\qquad
b(\vartheta)
\]
einen nichttrivialen gemeinsamen Teiler in $F[\vartheta]$ haben und
wir schreiben
\[
v_k(\vartheta)
=
\operatorname{ggT}(a(\vartheta)-c_k\cdot Db(\vartheta), b(\vartheta))
\in
F[\vartheta]
\]
für den grössten gemeinsamen Teiler.
Wir vermuten dann, dass die Summe der $c_k\log v_k(\vartheta)$ die
Stammfunktion
\[
\int
\frac{a(\vartheta)}{b(\vartheta)}
\overset{?}{=}
\sum_{k} c_k \log v_k(\vartheta)
\]
ist oder
\[
\frac{a(\vartheta)}{b(\vartheta)}
\overset{?}{=}
\sum_{k} D\biggl(c_k \log v_k(\vartheta)\biggr)
=
\sum_{k}\biggl(
Dc_k \log v_k(\vartheta)
+
c_k \frac{Dv_k(\vartheta)}{v_k(\vartheta)}
\biggr).
\]
Wenn dies tatsächlich die Stammfunktion in der Form sein soll, die
das Liouville-Prinzip verspricht, dann dürfen die Logarithmus-Terme
nicht vorkommen, d.~h.~es muss $Dc_k=0$ sein oder $c_k$ muss konstant
sein.

\begin{satz}[Rothstein-Trager]
\label{buch:intalg:risch:satz:rothstein-trager-allgemein}
\index{Satz!von Rothstein-Trager}
Gegeben ist eine logarithmische oder exponentielle Körpererweiterung
$F\subset F(\vartheta)$ und ein rationaler Teil wie in 
\eqref{buch:intalg:risch:eqn:a(theta)/b(theta)}.
\begin{enumerate}
\item Wenn die Stammfunktion von $a(\vartheta)/b(\vartheta)$ elementar ist,
dann sind die einzigen Elemente $c\in F$, für die
\[
a(\vartheta)-c\cdot Db(\vartheta)
\qquad\text{und}\qquad
b(\vartheta)
\]
einen nichttrivialen gemeinsamen Teiler haben, Konstanten.
\item
Falls die einzigen Elemente $c\in F$, für die 
$a(\vartheta)-c\cdot Db(\vartheta)$
und
$b(\vartheta)$
einen nichttrivialen gemeinsamen Teiler in $F[\vartheta]$ haben,
Konstanten sind, dann sind die Polynome
\begin{equation}
v_k(\vartheta)
=
\operatorname{ggT}(
a(\vartheta) - c_k\cdot Db(\vartheta), b(\vartheta)
).
\label{buch:intalg:risch:rothsteintrager:eqn:vk}
\end{equation}
teilerfremd und haben das Produkt
\[
b(\vartheta) = \prod_{k=1}^m v_k(\vartheta).
\]
\end{enumerate}
\end{satz}

Man beachte, dass dieser Satz noch nicht die ganze Aussage des
Satzes~\ref{buch:ratint:rothsteintrager:satz:rothsteintrager} auf den
Fall einer Erweiterung $F\subset F(\vartheta)$ verallgemeinert.
Es fehlt noch die Aussage, dass sich aus den Elementen $v_k(\vartheta)$
und Konstanten $c_k$ tatsächlich die Stammfunktion konstruieren lässt.
Im Fall einer exponentiellen Erweiterung ist auch $v_0$ noch nicht definiert.
Der Beweis dieser Aussagen ist jedoch von der Art der Erweiterung abhängig
und wird daher später für jede Art von Erweiterung separat nachgeholt.

\begin{proof}
Wir orientieren uns für den Beweis an der Vorgehensweise im Beweis von
Satz~\ref{buch:ratint:rothsteintrager:satz:rothsteintrager} über die
Methode von Rothstein-Trager für rationale Funktionen.

Falls die Stammfunktion elementar ist, dann hat sie nach 
Satz~\ref{buch:intalg:liouville:polynome:satz:exponentiell}
die Form
\[
\int \frac{a(\vartheta)}{b(\vartheta)}
=
v_0(\vartheta)
+
\sum_{k=1}^m c_k\log v_k(\vartheta)
\]
mit $v_0(\vartheta)\in F(\vartheta)$ und $v_k(\vartheta)\in F[\vartheta]$.
Nach Lemma~\ref{buch:intalg:risch:lemma:v0log} und
\ref{buch:intalg:risch:lemma:v0exp}
muss $v_0(\vartheta)=0$ bzw.~$v_0(\vartheta)\in F$ sein.

Seien jetzt die $c_k$ die einzigen Elemente $c\in F$, für die
$a(\vartheta)-c\cdot Db(\vartheta)$ und $b(\vartheta)$ einen nichttrivialen
gemeinsamen Teiler in $F[\vartheta]$ haben, und $v_k(\vartheta)$ wie in
\eqref{buch:intalg:risch:rothsteintrager:eqn:vk} definiert.

Ist $w=\operatorname{ggT}(v_i(\vartheta),v_j(\vartheta))$ mit $i\ne j$, dann
ist 
\begin{align*}
w&\mid (a(\vartheta)-c_i\cdot Db(\vartheta))\\
w&\mid (a(\vartheta)-c_j\cdot Db(\vartheta))\\
w&\mid \phantom{(}b(\vartheta).\phantom{)}
\end{align*}
Folglich ist $w$ auch ein Teiler von $(c_i-c_j) Db(\vartheta)\ne 0$.
Da $b(\vartheta)$ quadratfrei ist, können $b(\vartheta)$ und $Db(\vartheta)$
keinen gemeinsamen Teiler haben, also muss $w=1$ sein.
Die Polynome $v_k(\vartheta)$ sind also teilerfremd.

Wir betrachten jetzt das Produkt
\[
r(\vartheta) = \prod_{k=1}^m v_k(\vartheta).
\]
Da $v_k(\vartheta)\mid b(\vartheta)$ ist, gibt es ein Polynom $s(\vartheta)$
derart, dass
\begin{equation}
b(\vartheta) = r(\vartheta) s(\vartheta).
\label{buch:intalg:risch:rt:eqn:b=rs}
\end{equation}
Wir wollen zeigen, dass $s(\vartheta)=1$ ist.
Falls der Grad $\deg s(\vartheta)>0$ ist, suchen wir Elemente $z_0$, für
die 
$a(\vartheta) - z\cdot Db(\vartheta)$ und $s(\vartheta)$ einen nichttrivialen
gemeinsamen Teiler $t(\vartheta)$ haben.
Ein solcher Teiler müsste auch ein Teiler von $b(\vartheta)$ sein.
Die einzigen Zahlen $z$, für die dies möglich ist, müssten die Zahlen
$c_k$ sein, $z_0$ ist also eines der $c_k$.
Dann ist aber $t(\vartheta)$ ein Teiler von $v_k(\vartheta)$ und von
$s(\vartheta)$, er kommt somit im Produkt 
\eqref{buch:intalg:risch:rt:eqn:b=rs}
zweimal vor, im Widerspruch zur Tatsache, dass $b(\vartheta)$ quadratfrei
ist.
Der Widerspruch zeigt, dass $\deg s(\vartheta)=0$ ist und dass
$b(\vartheta)=r(\vartheta)$.
\end{proof}

Die von Satz~\ref{buch:intalg:risch:satz:rothstein-trager-allgemein}
formulierte Bedingung zur Bestimmung der Koeffizienten ist noch nicht
optimal.
Die Durchführung des euklidischen Algorithmus liefert aber tatsächlich
algebraische Gleichungen zur Bestimmung von $c_k$, wie spätere Beispiele
zeigen werden.
In Kapitel~\ref{chapter:resultante} wird mit der Resultanten ausserdem
ein elegantes algebraisches Werkzeug zur Bestimmung der Koeffizienten
angegeben.


%
% Hindernisse zur Existenz einer elementaren Stammfunktion
%
\subsubsection{Hindernisse zur Existenz einer elementaren Stammfunktion}
Die Skizze des Algorithmus im vorangegangenen Abschnitt hat gezeigt,
dass die Erweiterung der Vorgehensweise von Kapitel~\ref{chapter:ratint}
Hindernisse für die Existenz einer Stammfunktion finden wird.
Die Forderung, dass die Methode von Rothstein-Trager Konstanten
finden muss, obwohl die Koeffizienten der beteiligten Polynome
nichtkonstante Funktionen sein können, ist ein solches Hindernis.

Der polynomielle Teil wird weitere Hindernisse zutage fördern.
Ob die Stammfunktion des polynomiellen Teils existiert, hängt von
den Koeffizienten von $s(\vartheta)$ ab.
